\chapter{حجية الإجماع والقياس وإبطال الاستحسان}

\section{مقدمة في حجية الإجماع}
الحمد لله والصلاة والسلام على رسول الله، أما بعد؛ فقد عقد الإمام الشافعي رحمه الله تعالى باباً نفيساً في "الإجماع" ليبين حجية ما اجتمع عليه المسلمون مما ليس فيه نص حكم يُحكى عن النبي ﷺ.
فبين أن علماء الأمة إذا أجمعوا على مسألة وذكروا دليلها، فهذا هو الدليل. وأما إذا أجمعوا ولم يذكروا دليلاً نصياً، فإننا نتبعهم في إجماعهم؛ لأن سنة رسول الله ﷺ إن غابت عن بعضهم فلن تغيب عن جميعهم، ولأن عامة المسلمين وعلمائهم لا تجتمع على خلاف لسنة رسول الله ﷺ، ولا تجتمع على خطأ وضلالة أبداً.

\section{الأدلة على حجية الإجماع ولزوم الجماعة}
واستدل الشافعي رحمه الله على حجية الإجماع بحديث ابن مسعود رضي الله عنه أن النبي ﷺ قال: «نَضَّرَ الله عبداً سَمِعَ مَقَالَتي فَحَفِظَهَا وَوَعَاهَا وَأَدَّاهَا... ثَلَاثٌ لا يُغِلُّ عَلَيْهِنَّ قَلْبُ مُسْلِمٍ: إِخْلَاصُ العَمَلِ للهِ، وَالنَّصِيحَةُ لِلْمُسْلِمِينَ، وَلُزُومُ جَمَاعَتِهِمْ، فَإِنَّ دَعْوَتَهُمْ تُحِيطُ مِنْ وَرَائِهِمْ».
وكذلك بحديث عمر بن الخطاب رضي الله عنه بالجابية حين خطب الناس فقال: إن رسول الله ﷺ قام فينا مقامي فيكم، فقال: «...عَلَيْكُمْ بِالْجَمَاعَةِ، وَإِيَّاكُمْ وَالْفُرْقَةَ، فَإِنَّ الشَّيْطَانَ مَعَ الْوَاحِدِ وَهُوَ مِنَ الاثْنَيْنِ أَبْعَدُ، مَنْ أَرَادَ بُحْبُوحَةَ الْجَنَّةِ فَلْيَلْزَمِ الْجَمَاعَةَ».

\subsection{المعنى الصحيح للزوم الجماعة}
سُئل الشافعي: ما معنى أمر النبي ﷺ بلزوم جماعتهم؟
فقال: لا معنى له إلا لزوم إجماعهم؛ لأن أبدان المسلمين متفرقة في البلدان ولا يمكن لزومها، ولأن اجتماع الأبدان قد يضم المسلم والكافر والبر والفاجر. فلم يكن للزوم أبدانهم معنى، وإنما المعنى: لزوم ما هم عليه من التحليل والتحريم والطاعة. فمن قال بما تقول به جماعة المسلمين (في العقيدة والمنهج والأحكام) فقد لزم جماعتهم، ومن خالف ما أجمعوا عليه فقد خالف الجماعة.
وهذا الإجماع هو إجماع علماء السلف، ولا عبرة بدعاوى الإجماع المعاصرة التي تخالف ما كان عليه السلف الصالح، كمن يدعي الإجماع على النظم الوضعية والديمقراطية المستوردة المصادمة لحاكمية الشريعة.

\section{الاجتهاد والقياس: اسمان لمعنى واحد}
ثم انتقل الشافعي لبيان القياس، وذكر أن القياس والاجتهاد هما اسمان لمعنى واحد. وجماعهما: أن كل ما نزل بمسلم ففيه حكم لازم، إما أن يكون نصاً بعينه فيُتبع، وإما أن تطلب فيه الدلالة على سبيل الحق بالاجتهاد، وهذا الاجتهاد هو "القياس".

\subsection{التكليف بالظاهر دون الباطن}
وبيّن الشافعي أن العباد كُلفوا بإصابة الحق في الظاهر والباطن معاً فيما فيه نص، وكُلفوا بإصابة الحق في الظاهر (وهو ما ظهر لهم بالاجتهاد) فيما لا نص فيه.
وضرب لذلك أمثلة:
\begin{itemize}
    \item \textbf{استقبال القبلة:} من كان معاينًا للكعبة فهو يصيب الحق ظاهراً وباطناً بيقين. ومن غاب عنها فهو مأمور بالاجتهاد بالعلامات والنجوم، فإذا اجتهد وتوجه فقد أدى ما عليه، وهو مصيب في الظاهر، ولا يعلم الباطن (حقيقة الإصابة الدقيقة) إلا الله.
    \item \textbf{الحكم بعدالة الشهود:} نحن مأمورون بقبول شهادة العدل. ولا يُعرف العدل إلا بما يظهر لنا من صلاحه وملازمته للتقوى والمروءة، فنحكم بعدالته ظاهراً، وقد يكون في الباطن بخلاف ذلك، لكننا لم نُكلف بعلم الغيب.
    \item \textbf{الاجتهاد في جزاء الصيد:} في قوله تعالى ﴿فَجَزَاءٌ مِّثْلُ مَا قَتَلَ مِنَ النَّعَمِ﴾، يحكم العدلان بأقرب الأشياء شبهاً بالصيد المقتول من النعم، وهذا التقريب هو من باب الاجتهاد والقياس.
\end{itemize}

\subsection{أجر المجتهد}
واستدل بحديث عمرو بن العاص رضي الله عنه: «إذا حكم الحاكم فاجتهد ثم أصاب فله أجران، وإذا حكم فاجتهد ثم أخطأ فله أجر». فالمجتهد مأجور على بذل جهده لإعمال النصوص، فإن وافق الحق عند الله (الباطن) نال أجرين، وإن وافق الظاهر عنده وأخطأ الباطن نال أجراً واحداً.

\section{إبطال الاستحسان}
حذر الشافعي تحذيراً شديداً من القول بالاستحسان المجرد عن الدليل والقياس، وقرر أن أحداً لا يجوز له أن يفتي في دين الله إلا من جهة العلم (الكتاب، السنة، الإجماع، القياس).
فلا يجوز أن يُقال بالاستحسان (وهو القول بالرأي والتشهي دون مستند شرعي). وضرب لذلك مثلاً دنيوياً: إذا أتلف شخص سلعة لآخر، فلا يُلجأ في تقييمها إلا لخبير خابر بالسوق يقيسها بمثيلاتها، ولا يجوز له أن يلقي سعراً بالتعسف والاستحسان. فإذا كان هذا في أموال الدنيا، فكيف بحلال الله وحرامه؟
فالاستحسان المجرد ما هو إلا تلذذ واتباع للهوى، والواجب على العالِم ألا يقول إلا من جهة العلم، ولا يحل له أن يتخبط في دين الله برأيه المجرد.